% !TEX program = lualatex
% 10: 情報セキュリティ基礎

\documentclass[aspectratio=43,professionalfonts,handout,10pt]{beamer}
\usepackage{beamer_template}

\newcommand{\LectureNum}{10}
\newcommand{\LectureTitle}{情報セキュリティ基礎}
\newcommand{\CourseName}{情報リテラシー}
\newcommand{\TextPages}{08-09}

\begin{document}
% --- Slide 1: title ---

\begin{frame}{\CourseName\quad 第\LectureNum 回「\LectureTitle 」}
  {\small \TermName\quad \TeacherName\quad ／\quad 教科書 \TextPages}

  \vfill

  \begin{block}{今日の流れ}
    \begin{enumerate}
      \item 不正アクセスとパスワードの不正取得
      \item 情報セキュリティとは（CIA）
      \item 個人や組織の取り組み
      \item 法律と制度・情報モラルと個人の責任
    \end{enumerate}
  \end{block}
  \vfill

  \begin{block}{今日の目標}
    \begin{itemize}
      \item 情報セキュリティの必要性を理解し，基本的な対策について説明できる。
    \end{itemize}
  \end{block}
  \vfill

  \begin{exampleblock}{キーワード}
    \small \textbf{不正アクセス} ／ \textbf{マルウェア} ／ \textbf{ソーシャルエンジニアリング} ／ \textbf{フィッシング} ／ \textbf{CIA}（機密性・完全性・可用性）／ \textbf{情報モラル} ／ \textbf{デジタルタトゥー}
  \end{exampleblock}
\end{frame}

% --- Slide 2: 不正アクセス ---

\begin{frame}{不正アクセス（教科書 p.18）}
  \begin{block}{不正アクセスとは}
    ネットワークを利用して，\textbf{正当な権限のないコンピュータにアクセスする行為}
    \begin{itemize}
      \item 個人情報の盗取・データの改変などの被害が起こる
      \item \textbf{不正アクセス禁止法}で規制されている
    \end{itemize}
  \end{block}

  \vfill

  \begin{alertblock}{注意}
    他人のID・パスワードを無断で使ってログインするだけで違法。\\
    「試しにやってみた」も対象になる。未成年にも適用される。
  \end{alertblock}

  \begin{slidenote}
    不正アクセス禁止法は1999年公布。アクセスの結果に関わらず，不正アクセス行為自体が違法になる。
  \end{slidenote}
\end{frame}

% --- Slide 3: パスワードの不正取得 ---

\begin{frame}{パスワードの不正取得（教科書 p.18）}
  {\small 認証（ユーザID・パスワード）を不正に取得する主な手法：}

  \begin{columns}[T]
    \begin{column}{0.48\textwidth}
      \begin{exampleblock}{マルウェア}
        \small
        コンピュータに害を与えるソフトウェアの総称（ウイルス・ランサムウェアなど）。\\
        感染するとパスワード盗取・データ破壊が起こる。
      \end{exampleblock}
    \end{column}
    \begin{column}{0.48\textwidth}
      \begin{exampleblock}{ソーシャルエンジニアリング}
        \small
        人の心理的弱点を利用して情報を盗む手口。
        \begin{itemize}
          \item \textbf{フィッシング}：偽メール・偽サイトでID・パスワードを入力させる
          \item なりすまし電話・ショルダーハック
        \end{itemize}
      \end{exampleblock}
    \end{column}
  \end{columns}

  \begin{exampleblock}{演習}
    フィッシング詐欺に引っかかりやすい状況はどんなときか，話し合ってみよう。
  \end{exampleblock}

  \begin{slidenote}
    身近な例：銀行やAmazonを装った偽メールが届き，偽サイトでログイン情報が盗まれる。URLや送信元アドレスを必ず確認しよう。
  \end{slidenote}
\end{frame}

% --- Slide 4: 情報セキュリティとは（CIA） ---

\begin{frame}{情報セキュリティとは}
  \begin{block}{情報セキュリティの定義}
    情報の\textbf{機密性・完全性・可用性}を維持すること（CIA三原則）
  \end{block}

  \vfill

  \begin{columns}[T]
    \begin{column}{0.32\textwidth}
      \begin{exampleblock}{機密性\\\small Confidentiality}
        \small
        許可された人だけが情報にアクセスできる状態を保つ
      \end{exampleblock}
    \end{column}
    \begin{column}{0.32\textwidth}
      \begin{exampleblock}{完全性\\\small Integrity}
        \small
        情報が正確で改ざんされていない状態を保つ
      \end{exampleblock}
    \end{column}
    \begin{column}{0.32\textwidth}
      \begin{exampleblock}{可用性\\\small Availability}
        \small
        必要なときに情報を利用できる状態を保つ
      \end{exampleblock}
    \end{column}
  \end{columns}

  {\small \textbf{主な脅威：}不正アクセス・マルウェア・災害・操作ミス・内部不正}

  \begin{slidenote}
    不正アクセスは主に機密性を脅かす。マルウェアはウイルス（完全性）・ランサムウェア（可用性）・スパイウェア（機密性）など，CIAのすべてに影響しうる。
  \end{slidenote}
\end{frame}

% --- Slide 5: 個人や組織の取り組み ---

\begin{frame}{個人や組織の取り組み（教科書 p.19）}
  \resizebox{\textwidth}{!}{%
  \begin{tabular}{lll}
    \toprule
    対策の種類 & 内容 & 具体例 \\
    \midrule
    技術的対策 & システムや技術による防御 & ファイアウォール・暗号化・ウイルス対策ソフト \\
    物理的対策 & 物理的な環境の管理 & 入退室の管理・施錠管理 \\
    人的対策 & 人による適切な運用 & セキュリティ教育・パスワード管理 \\
    \bottomrule
  \end{tabular}}

  \vfill

  \begin{block}{セキュリティポリシー（教科書 p.19）}
    \small 組織における情報セキュリティの\textbf{方針・規則をまとめた文書}。全構成員が遵守する。
  \end{block}

  \begin{exampleblock}{演習}
    自分が使っているサービス（SNS・学内システム等）でどんなセキュリティ対策が施されているか，話し合ってみよう。
  \end{exampleblock}

  \begin{slidenote}
    技術・物理・人的の3種類をバランスよく組み合わせることが重要。どんなに強固なシステムでも，人がパスワードを教えてしまえば意味がない。
  \end{slidenote}
\end{frame}

% --- Slide 6: 法律と制度 ---

\begin{frame}{法律と制度（教科書 p.20）}
  \begin{block}{情報社会に関わる主な法律}
    \small
    \begin{tabular}{ll}
      \toprule
      年 & 主な法律・社会の動き \\
      \midrule
      1999年 & 不正アクセス禁止法 公布 \\
      2001年 & プロバイダ責任制限法 公布 \\
      2003年 & 個人情報保護法 公布 \\
      2022年 & 侮辱罪の厳罰化・法改正 \\
      \bottomrule
    \end{tabular}
  \end{block}

  \vfill

  \begin{alertblock}{ポイント}
    法律・制度だけでは不十分。情報社会の変化に対応しつづける姿勢と，\textbf{個人が適切に判断・行動する力}が必要
  \end{alertblock}

  \begin{slidenote}
    法律は技術の進化に後追いになりやすい。「法律がないからOK」ではなく，情報モラルを持って行動することが大切。
  \end{slidenote}
\end{frame}

% --- Slide 7: 情報モラルと個人の責任 ---

\begin{frame}{情報モラルと個人の責任（教科書 p.21）}
  \begin{columns}[T]
    \begin{column}{0.48\textwidth}
      \begin{block}{情報リテラシー}
        \small
        必要な情報を適切に\textbf{収集・分析・判断・表現}できる能力。\\
        特に\textbf{情報の真偽を適切に判断}することが重要。
        \begin{itemize}
          \item 著作権・個人情報に配慮する
          \item ネット炎上などのリスクを考慮する
        \end{itemize}
      \end{block}
    \end{column}
    \begin{column}{0.48\textwidth}
      \begin{alertblock}{デジタルタトゥー}
        \small
        \begin{itemize}
          \item ネット上に一度公開した情報は完全には消せない
          \item 投稿・写真・発言が長期にわたって残り続ける
          \item 将来の就職・人間関係に影響を与える
        \end{itemize}
      \end{alertblock}
    \end{column}
  \end{columns}

  \begin{exampleblock}{演習}
    SNSに投稿する前に確認すべきことを3つ以上挙げ，話し合ってみよう。
  \end{exampleblock}

  \begin{slidenote}
    情報モラルとは，技術スキルだけでなく，判断力・倫理観の両方が必要な力。デジタルタトゥーは「消した」つもりでもスクリーンショットやキャッシュで残る。
  \end{slidenote}
\end{frame}

% --- チェックリスト ---

\begin{frame}{まとめ・チェックリスト}
  \begin{block}{自己チェック（できたら $\checkmark$ をつけよう）}
    \begin{itemize}
      \item[$\square$] 不正アクセス・マルウェア・ソーシャルエンジニアリングの脅威を理解した
      \item[$\square$] 情報セキュリティのCIA三原則（機密性・完全性・可用性）を説明できる
      \item[$\square$] 技術的・物理的・人的の3種類のセキュリティ対策を学んだ
      \item[$\square$] 情報モラルの重要性と，情報社会に関わる法律の変化を理解した
      \item[$\square$] デジタルタトゥーの意味と，ネット上の行動に伴う責任を理解した
    \end{itemize}
  \end{block}

  \vfill

  {\small 質問があれば授業後またはTeamsで受け付けます。}
\end{frame}

\end{document}
